\mychapter{Continuidade e Validação Final}{cap:plano_continuidade}\lhead{CONTINUIDADE E VALIDAÇÃO FINAL}

A metodologia AIPE está formulada e possui duas aplicações experimentais. A Bahia e o piloto M5 têm matrizes completas. A avaliação unificada amplia os seletores na Bahia e organiza a continuidade do M5. Sua consolidação precede a avaliação final; a rodada parcial não é apresentada como experimento concluído. O piloto já orientou diagnóstico e recuo; integra o conjunto de desenvolvimento. A avaliação nacional está fora do escopo experimental desta versão.

\section{Separar desenvolvimento e confirmação}
\label{sec:estudos_caso_adicionais}
A próxima amostra deve oferecer mais exemplos por classe e registrar suporte de cada perfil. A concentração natural do M5 em séries esparsas impede presumir cobertura equilibrada. Elegibilidade e estratificação usam somente o treino; séries reservadas não podem orientar o ajuste do seletor.

O pipeline unificado usa cinco grupos internos por loja para escolher a margem de recuo, mantendo a loja externa excluída em cada origem. Outros hiperparâmetros e eventual simplificação de rótulos devem ser definidos somente no desenvolvimento. O teste final fica isolado. Uma melhora encontrada no piloto não é uma segunda replicação independente \citep{CawleyTalbot2010}. Ajustes de políticas compatíveis podem ser reaproveitados após auditoria de dados, configuração e código; mudanças científicas exigem novo ajuste da candidata afetada.

\section{Harmonizar o problema de decisão}
A avaliação unificada aplica pesos zero de ruptura e amplificação à calibração da CappedBaseStock nas duas bases. As tabelas desta versão ainda refletem as matrizes anteriores; a atualização depende da conferência das novas saídas. Também é necessário comparar restrições de serviço com escopo claro: série, perfil ou rede. A meta de 0,70 e os custos atuais são parâmetros experimentais, não exigências empresariais já justificadas.

A comparação reduzida por consequências e recuo fornece um candidato para continuidade, com benefício concentrado em três intervenções. A instabilidade dos rótulos também motiva examinar conjuntos de candidatas quase empatadas e alvos sensíveis ao custo de erro. Essa análise deve ocorrer no desenvolvimento, sem selecionar a formulação pelo teste reservado. GA--DQN continua acompanhado por GA e DQN, com o esforço adicional informado.

\section{Entrega gerencial e avaliação sequencial}
A próxima etapa de desenvolvimento deve completar o contrato de recomendação: origem e qualidade das entradas, família e parâmetros finais, alternativas, consequências esperadas, hipóteses e ligação com a avaliação posterior. A prévia da Bahia reaproveita indicadores persistidos; parâmetros e trajetórias ausentes não podem ser inferidos dos totais. Sua recuperação requer execução adicional quando não houver artefato exato, preservando a matriz de origem. O julgamento humano e a efetividade do apoio exigem estudo com usuários ou em campo.

A cadência informada para a empresa-fonte é uma revisão por ciclo comercial. Uma avaliação dessa rotina deve fixar período de ação, agregação, revisão do PSE, horizonte, janela e atualização dos modelos. Cada estratégia conserva seu próprio estoque e pedidos em trânsito; mudanças de família não reinicializam o sistema. O calendário físico e as convenções de custo precisam permanecer comparáveis ao testar outras cadências. A execução por blocos continua válida em seu escopo e não deve ser modificada silenciosamente para representar esse novo desenho.

\section{Sensibilidades e eficiência}
\label{sec:robustez_pooling}
Prazo, custos relativos, estoque inicial, tratamento terminal e horizonte podem alterar o ranking e a viabilidade física. Seu efeito deverá ser separado do efeito do seletor. Calibração agrupada por perfil é outra hipótese: não equivale a seleção de família por perfil e exige comparação própria.

A escolha da unidade central de processamento (\textit{central processing unit}, CPU) ou da unidade de processamento gráfico (\textit{graphics processing unit}, GPU) deve considerar tempo total sob a carga de uso prevista. O piloto utilizou oito processos CPU; essa configuração operacional não é uma vantagem científica de uma política. A eficiência computacional deve ser relatada como propriedade da execução, separada da qualidade das recomendações. A relevância prática será discutida por magnitude, incerteza, deteriorações e contexto, sem impor economia percentual universal. Inferência de não inferioridade do atendimento, se pretendida, exigirá margem justificada e desenho específico.

\section{Critérios de fechamento}
\begin{table}[htbp]\centering\small
\caption{Critérios de conclusão, sem datas ainda não acordadas.}\label{tab:cronograma}
\begin{tabular}{@{}p{4cm}p{8.8cm}@{}}\toprule
Etapa & Evidência necessária\\\midrule
Formulação do AIPE & Entradas, etapas, saídas e limites explicitados independentemente das bases\\
Protocolo final & Objetivo, suporte e calibrações fixados antes do teste\\
PSE & Comparação independente com global, perfil e controles simples\\
Robustez & Falhas físicas, portfólio, seleção e sensibilidades separadas\\
Redação & Afirmações alinhadas à evidência e revisão dos autores\\\bottomrule
\end{tabular}\end{table}

O artigo concentra a arquitetura e sua avaliação operacional; a dissertação mantém fundamentos, formulação e limites mais amplos. Autoria, declarações, disponibilidade de dados e elementos institucionais precisam ser fechados pelos responsáveis antes de submissão ou depósito.
